% Options for packages loaded elsewhere
\PassOptionsToPackage{unicode}{hyperref}
\PassOptionsToPackage{hyphens}{url}
%
\documentclass[
]{article}
\usepackage{amsmath,amssymb}
\usepackage{iftex}
\ifPDFTeX
  \usepackage[T1]{fontenc}
  \usepackage[utf8]{inputenc}
  \usepackage{textcomp} % provide euro and other symbols
\else % if luatex or xetex
  \usepackage{unicode-math} % this also loads fontspec
  \defaultfontfeatures{Scale=MatchLowercase}
  \defaultfontfeatures[\rmfamily]{Ligatures=TeX,Scale=1}
\fi
\usepackage{lmodern}
\ifPDFTeX\else
  % xetex/luatex font selection
\fi
% Use upquote if available, for straight quotes in verbatim environments
\IfFileExists{upquote.sty}{\usepackage{upquote}}{}
\IfFileExists{microtype.sty}{% use microtype if available
  \usepackage[]{microtype}
  \UseMicrotypeSet[protrusion]{basicmath} % disable protrusion for tt fonts
}{}
\makeatletter
\@ifundefined{KOMAClassName}{% if non-KOMA class
  \IfFileExists{parskip.sty}{%
    \usepackage{parskip}
  }{% else
    \setlength{\parindent}{0pt}
    \setlength{\parskip}{6pt plus 2pt minus 1pt}}
}{% if KOMA class
  \KOMAoptions{parskip=half}}
\makeatother
\usepackage{xcolor}
\usepackage[margin=1in]{geometry}
\usepackage{graphicx}
\makeatletter
\def\maxwidth{\ifdim\Gin@nat@width>\linewidth\linewidth\else\Gin@nat@width\fi}
\def\maxheight{\ifdim\Gin@nat@height>\textheight\textheight\else\Gin@nat@height\fi}
\makeatother
% Scale images if necessary, so that they will not overflow the page
% margins by default, and it is still possible to overwrite the defaults
% using explicit options in \includegraphics[width, height, ...]{}
\setkeys{Gin}{width=\maxwidth,height=\maxheight,keepaspectratio}
% Set default figure placement to htbp
\makeatletter
\def\fps@figure{htbp}
\makeatother
\setlength{\emergencystretch}{3em} % prevent overfull lines
\providecommand{\tightlist}{%
  \setlength{\itemsep}{0pt}\setlength{\parskip}{0pt}}
\setcounter{secnumdepth}{-\maxdimen} % remove section numbering
\usepackage{booktabs}
\usepackage{longtable}
\usepackage{array}
\usepackage{multirow}
\usepackage{wrapfig}
\usepackage{float}
\usepackage{colortbl}
\usepackage{pdflscape}
\usepackage{tabu}
\usepackage{threeparttable}
\usepackage{threeparttablex}
\usepackage[normalem]{ulem}
\usepackage{makecell}
\usepackage{xcolor}
\ifLuaTeX
  \usepackage{selnolig}  % disable illegal ligatures
\fi
\usepackage{bookmark}
\IfFileExists{xurl.sty}{\usepackage{xurl}}{} % add URL line breaks if available
\urlstyle{same}
\hypersetup{
  pdftitle={Team 1 - Lab 1},
  pdfauthor={Nicole, Brian, Anna},
  hidelinks,
  pdfcreator={LaTeX via pandoc}}

\title{Team 1 - Lab 1}
\usepackage{etoolbox}
\makeatletter
\providecommand{\subtitle}[1]{% add subtitle to \maketitle
  \apptocmd{\@title}{\par {\large #1 \par}}{}{}
}
\makeatother
\subtitle{Forecasting with ARIMA and ETS models}
\author{Nicole, Brian, Anna}
\date{April 21, 2025}

\begin{document}
\maketitle

{
\setcounter{tocdepth}{2}
\tableofcontents
}
\section{Data}\label{data}

Discuss what parts of the data sets you will work with.

\section{Question your team will
address}\label{question-your-team-will-address}

What model structure best fits the different states and regions (Total
and Pacific Northwest) represented in the yearly data set? How does
forecasting accuracy compare across these regions? How does model
performance change when you drop years? This failed so start dropping
months - address Brian's real issue of moving from a seasonal model to a
non seasonal (CA) because of data dropouts) Practical Application: You
might want to know when you can get the most value out of your catch -
how does the value per ton change over time?

\section{Method you will use}\label{method-you-will-use}

\subsection{Initial plan}\label{initial-plan}

\subsection{What you actually did}\label{what-you-actually-did}

Example, ``We were able to do our plan fairly quickly, so after writing
the basic code. We compared the accuracy of forecasts for different time
periods using 20-years for training and 5-year forecasts each time. We
compared the RMSE, MAE, and MAPE accuracy measures.''

\section{Diagnostics and preliminary
exploration}\label{diagnostics-and-preliminary-exploration}

To address question 1, we subset the data for each west coast state,
Pacific Northwest which includes WA, OR, and CA, and the total of all
states combined, into a training data set and test data set. The
training data sets spanned the years 1950-2005, and the test data set
was 2006-2016. When we plotted the data (Figure 1), there was no obvious
problems with stationarity. There might be a slight downwards trend, and
there are a few gaps in the California data, but it is not immediately
obvious if the data is non-stationary.

\includegraphics{Lab-1-Team-1_files/figure-latex/unnamed-chunk-2-1.pdf}

\subsection{Use ACF and PACF}\label{use-acf-and-pacf}

\includegraphics{Lab-1-Team-1_files/figure-latex/unnamed-chunk-3-1.pdf}
\includegraphics{Lab-1-Team-1_files/figure-latex/unnamed-chunk-3-2.pdf}
\includegraphics{Lab-1-Team-1_files/figure-latex/unnamed-chunk-3-3.pdf}

What did you learn? Also try decomposition. \emph{I might need help
interpreting these plots}

\subsection{Test for stationarity}\label{test-for-stationarity}

When we run ndiffs on the states' data, Oregon, Washington, the Pacific
Northwest region determined one first-difference is required to make the
data stationary.

\begin{verbatim}
##      Alaksa California Oregon Washington Total PNW
## [1,]      0          0      1          1     1   1
\end{verbatim}

\section{Results}\label{results}

We fit multiple ARIMA and Holt models to the yearly time series data for
each state, the PNW region (WA, OR, CA), and the total chinook landings
for all states combined. ARIMA models where p, d, or q were greater than
1 had significantly higher AIC values than more parsimonious ARIMA
models and are excluded from this write-up. We then forecasted 10 years
of data using the six ARIMA models, tested the forecasts against the
test data for each region. Using root mean squared error, we selected
the best fit model for each state or region. The Holt model was the best
fit for Alaska, California, and the Pacific Northwest region. An
ARIMA(1,0,0) model was the best fit for Washington and Oregon. An
ARIMA(0,1,0) model was the best fit for the entire Pacific Coast chinook
landing data (log metric tons). See the figures below for the six
regional time series and forecasts from the best fit model for each. We
will also provide a brief interpretation of the best fit models, and
check of the residuals using a Ljung-Box test.

\begin{verbatim}
## Series: aktrain 
## ARIMA(1,0,0) with non-zero mean 
## 
## Coefficients:
##          ar1    mean
##       0.7069  8.5064
## s.e.  0.0909  0.0716
## 
## sigma^2 = 0.02773:  log likelihood = 21.6
## AIC=-37.2   AICc=-36.74   BIC=-31.12
\end{verbatim}

\begin{verbatim}
## Series: catrain 
## ARIMA(1,0,0) with non-zero mean 
## 
## Coefficients:
##          ar1    mean
##       0.4235  7.8462
## s.e.  0.1195  0.0850
## 
## sigma^2 = 0.1432:  log likelihood = -24.12
## AIC=54.25   AICc=54.71   BIC=60.32
\end{verbatim}

\begin{verbatim}
## Series: ortrain 
## ARIMA(0,1,0) 
## 
## sigma^2 = 0.2139:  log likelihood = -36.27
## AIC=74.55   AICc=74.62   BIC=76.57
\end{verbatim}

\begin{verbatim}
## Series: watrain 
## ARIMA(0,1,0) 
## 
## sigma^2 = 0.04634:  log likelihood = 6.43
## AIC=-10.87   AICc=-10.79   BIC=-8.86
\end{verbatim}

\begin{verbatim}
## Series: pnwtrain 
## ARIMA(0,1,0) 
## 
## sigma^2 = 0.08109:  log likelihood = -8.96
## AIC=19.91   AICc=19.99   BIC=21.92
\end{verbatim}

\begin{verbatim}
## Series: totaltrain 
## ARIMA(0,1,0) 
## 
## sigma^2 = 0.03084:  log likelihood = 17.63
## AIC=-33.25   AICc=-33.18   BIC=-31.25
\end{verbatim}

\paragraph{Alaska}\label{alaska}

The best fit model for the Alaska chinook landing data was a Holt model
with the following equation: \(y_t=l_{t-1}+b_{t-1}+\epsilon_t\)\\
\(l_t=l_{t-1}+b_{t-1}+0.9999*\epsilon_t\)\\
\(b_t =b_{t-1} + 0.0001*\epsilon_t\) This model indicates a slight
downward trend that remains fairly consistent over time (beta=0.0001),
and gives almost all weight to the most recent observations
(alpha=0.9999).

\begin{verbatim}
##          0,0,0     0,1,0     1,0,0     1,1,0    1,1,1      Holt
## RMSE 0.6611034 0.6579779 0.6641441 0.6567943 0.655056 0.6447745
\end{verbatim}

\includegraphics{Lab-1-Team-1_files/figure-latex/unnamed-chunk-6-1.pdf}

\begin{verbatim}
## 
##  Ljung-Box test
## 
## data:  Residuals from ETS(A,A,N)
## Q* = 5.4977, df = 10, p-value = 0.8556
## 
## Model df: 0.   Total lags used: 10
\end{verbatim}

\paragraph{California}\label{california}

The best fit model for California was also a Holt model with the
following equation: \(y_t=l_{t-1}+b_{t-1}+\epsilon_t\)
\(l_t=l_{t-1}+b_{t-1}+0.0005*\epsilon_t\)
\(b_t=b_{t-1}+0.0005*\epsilon_t\)

The model indicates a slightly declining trend, but small beta value
(beta=0.0005) suggest the level and trend are relatively stable over
time. The model also puts very little weight on recent observations
(alpha=0.0005).

\begin{verbatim}
##         0,0,0   0,1,0    1,0,0    1,1,0    1,1,1     Holt
## RMSE 3.053917 2.96879 3.048373 3.014439 3.005601 2.881865
\end{verbatim}

\includegraphics{Lab-1-Team-1_files/figure-latex/unnamed-chunk-7-1.pdf}

\begin{verbatim}
## 
##  Ljung-Box test
## 
## data:  Residuals from ETS(A,A,N)
## Q* = 37.425, df = 10, p-value = 4.776e-05
## 
## Model df: 0.   Total lags used: 10
\end{verbatim}

\paragraph{Oregon}\label{oregon}

The best fit model for Oregon was a first-order ARIMA model with no
differencing and no moving average. The time series is stationary and
auto-correlated around a mean of 7.5190, with observations being
moderately dependent on recent values (ar1\textless1).
\(y_t=0.7303*y_{t-1}+2.026+\epsilon_t\)
\(\epsilon_t~N(0,\sigma^2=0.1876)\)

\begin{verbatim}
##          0,0,0     0,1,0     1,0,0     1,1,0     1,1,1     Holt
## RMSE 0.8713649 0.6326566 0.5259991 0.5739301 0.5861011 0.737802
\end{verbatim}

\includegraphics{Lab-1-Team-1_files/figure-latex/unnamed-chunk-8-1.pdf}

\begin{verbatim}
## 
##  Ljung-Box test
## 
## data:  Residuals from ARIMA(1,0,0) with non-zero mean
## Q* = 7.3609, df = 9, p-value = 0.5996
## 
## Model df: 1.   Total lags used: 10
\end{verbatim}

\paragraph{Washington}\label{washington}

Like Oregon, the best fit model for Washington was a first-order ARIMA
model with no differencing and not moving average. The values differ
slightly, with Washington having a slightly higher average than Oregon
(intercept=7.9066), and autoregression component (ar1) of 0.8878,
meaning the series shows strong positive correlation.

\(y_t=0.8878*y_{t-1}+0.8871+\epsilon_t\)
\(\epsilon_t~N(0,\sigma^2=0.0436)\)

\begin{verbatim}
##          0,0,0     0,1,0     1,0,0     1,1,0     1,1,1      Holt
## RMSE 0.3707379 0.3193469 0.2504157 0.3188799 0.3192479 0.3744703
\end{verbatim}

\includegraphics{Lab-1-Team-1_files/figure-latex/unnamed-chunk-9-1.pdf}

\begin{verbatim}
## 
##  Ljung-Box test
## 
## data:  Residuals from ARIMA(1,0,0) with non-zero mean
## Q* = 5.206, df = 9, p-value = 0.816
## 
## Model df: 1.   Total lags used: 10
\end{verbatim}

\paragraph{PNW Region}\label{pnw-region}

The model that best fit the Pacific Northwest region as a whole was a
Holt model with the following equation:
\(y_t=l_{t-1}+b_{t-1}+\epsilon_t\)\\
\(l_t=l_{t-1}+b_{t-1}+0.9619*\epsilon_t\)\\
\(b_t =b_{t-1} + 0.0001*\epsilon_t\)

The small beta value suggests that the trend is stable over time, and
the alpha value being close to 1 indicates that the model assigns
greater weight to recent observations. The model has identified a
slightly negative trend in chinook landings over time (b=-0.0108).

\begin{verbatim}
##          0,0,0     0,1,0     1,0,0     1,1,0     1,1,1      Holt
## RMSE 0.8135519 0.6220047 0.7505696 0.6286976 0.6289971 0.5934286
\end{verbatim}

\includegraphics{Lab-1-Team-1_files/figure-latex/unnamed-chunk-10-1.pdf}

\begin{verbatim}
## 
##  Ljung-Box test
## 
## data:  Residuals from ETS(A,A,N)
## Q* = 16.47, df = 10, p-value = 0.08694
## 
## Model df: 0.   Total lags used: 10
\end{verbatim}

\paragraph{Total chinook landings}\label{total-chinook-landings}

The best fit model for total chinook landings was a first-order ARIMA
with first difference since the data was non-stationary, and no
seasonality: \(y_t = 0.8878*y_{t-1}+0.8871+\epsilon_t\)\\
\(\epsilon_t~N(0,\sigma^2=0.0436)\)\\

The alpha value suggests weak, negative autocorrelation, but it is not
statistically significant.

\begin{verbatim}
##         0,0,0   0,1,0     1,0,0     1,1,0     1,1,1      Holt
## RMSE 0.704873 0.56315 0.6618483 0.5659348 0.5694472 0.5312216
\end{verbatim}

\includegraphics{Lab-1-Team-1_files/figure-latex/unnamed-chunk-11-1.pdf}

\begin{verbatim}
## 
##  Ljung-Box test
## 
## data:  Residuals from ARIMA(1,1,0)
## Q* = 17.185, df = 9, p-value = 0.04589
## 
## Model df: 1.   Total lags used: 10
\end{verbatim}

\includegraphics{Lab-1-Team-1_files/figure-latex/unnamed-chunk-11-2.pdf}

In summary, none of these models contained a seasonality component. The
model structure differed slightly between regions, but none of the
models found significant changes in trend over time. Many of them
behaved functionally like a random walk.

\emph{Forecasting Accuracy Comparison}

\begin{verbatim}
##                 AK Model  CA Model  OR Model  WA Model PNW Model Total Model
## AK Forecast    0.6447745 0.4089038 0.7335605 0.3607922 0.7691459   1.4183031
## CA Forecast    3.5268146 2.8818654 2.7096443 2.9523386 3.6400794   4.1818441
## OR Forecast    1.7887483 0.9452960 0.5259991 1.0418512 1.9274951   2.5737358
## WA Forecast    0.8638981 0.3241475 0.5095812 0.2504157 1.0009427   1.6417700
## PNW Forecast   0.4860730 0.7278896 1.0296953 0.5795634 0.5934286   1.1656268
## Total Forecast 0.3771229 1.2245373 1.5626548 1.0840916 0.2934693   0.5659348
\end{verbatim}

\section{Discussion}\label{discussion}

In summary, none of these models contained a seasonality component. The
model structure differed slightly between regions, but none of the
models found significant changes in trend over time. Many of them
behaved functionally like a random walk.

When we compared forecasting accuracy between models fit to one region
against test data from another region, we found some interesting
results. In some cases, another state's forecasted data was a better fit
than the best fit model for that state. This was the case for Alaska,
whose test data was more accurately forecasted by the ARIMA model fit to
the Washington data. Only the Washington and Oregon model forecasts were
the best fit for their state's test data. This is the opposite of what I
would have expected, given our original hypothesis that Washington and
Oregon's time series look very similar. I would have expected forecasts
from these two states to not be accurate when tested against other
states' data, since California and Alaska's data looked different, but
the Washington and Oregon models performed fairly well across the board.
The large regional models (PNW and total) did not perform well, which
does support the idea that forecasts for salmon fisheries should be made
on smaller regional scales that are ecologically relevant to salmon
populations.

\section{Description of each team member's
contributions}\label{description-of-each-team-members-contributions}

Each member of the team took on one issue and completed the code and the
analysis. Halfway through, we consulted other team members for feedback
and adjust analysis based on the joint discussion. All team members
participated in the lab write up and review the results.

\end{document}
